% --------------------------------------------------------------------------
%		Document style
% --------------------------------------------------------------------------

\documentclass[11pt,a4paper,reqno,lualatex]{amsart} % reqnoは数式番号を右側表示(amsartのデフォルトは左)
%\documentclass[11pt,a4paper,oneside,lualatex]{ltjsarticle}
%\usepackage{luatexja} %日本語非対応のドキュメントクラスでも日本語を表示してくれる. 日本語のコメントを入れるときだけコメントアウトを外し、最終原稿では日本語を無くしてこの行を再びコメントアウトする
% SumatraPDFの逆順検索でエラーが出た時は以下のコマンドラインをSumatraPDFの設定→オプションで入力する
% "C:\Program Files\texstudio\texstudio.exe" "%f" -line %l

\usepackage[marginparwidth=0pt,margin=20truemm]{geometry} % margin
%\setcounter{tocdepth}{1} % 目次に出すのは\sectionまで

% --------------------------------------------------------------------------
%		Packages and commands
% --------------------------------------------------------------------------

\usepackage{mypackage} % My packages
\usepackage{mycommand} % My commands

% --------------------------------------------------------------------------
%		Hyperlinks
% --------------------------------------------------------------------------

\usepackage[luatex, pdfencoding=auto,hypertexnames=false]{hyperref}
\hypersetup{colorlinks=false}

% --------------------------------------------------------------------------
%		Theorem-like environments &  cross-reference
% --------------------------------------------------------------------------

\numberwithin{equation}{section} % Setting of equation numbers 
\newtheorem{theoremcounter}{}[section] % Setting of theorem numbers 
\usepackage{mytheoremeng} % My theorem environments (English) with cleveref package

% --------------------------------------------------------------------------
\begin{document}
% --------------------------------------------------------------------------


\title[Short title]{My template for works}
\author[Y. Murakami]{Yuya Murakami}
\address{Faculty of Mathematics, Kyushu University, 744, Motooka, Nishi-ku, Fukuoka, 819-0395, JAPAN}
\email{murakami.yuya.896@m.kyushu-u.ac.jp}
\thanks{The author is supported by...}

\date{\today}

\maketitle

% --------------------------------------------------------------------------

\begin{abstract}
	In this paper, we introduce...
	We show that...
\end{abstract}

% --------------------------------------------------------------------------

\tableofcontents

% --------------------------------------------------------------------------

\section{Introduction} \label{sec:intro}

% --------------------------------------------------------------------------

\begin{thm}[{\cite{test}}] \label{thm:main}
	hogehoge
\end{thm}

\begin{thm} \label{thm:main2}
	hogehoge
\end{thm}

\begin{thm} \label{thm:main3}
	hogehoge
\end{thm}

\begin{lem} \label{lem:1}
	hogehoge
\end{lem}

\begin{equation} \label{eq:1}
	abc
\end{equation}

\begin{equation} \label{eq:2}
	abc
\end{equation}

By \zcref{thm:main,thm:main2,thm:main3,lem:1,eq:1} , we have...

\begin{lem} \label{lem:2}
	\begin{enumerate}
		\item \label{item:lem:2:hoge} hogehoge
		\item \label{item:lem:2:huga} hugahuga
	\end{enumerate}
\end{lem}

By \zcref{lem:2} \zcref{item:lem:2:hoge,item:lem:2:huga}, we have...

\vspace{1em}

This paper will be organized as follows. 
In \zcref{sec:preliminaries}, we prepare...
In \zcref{sec:3}, we define...

% --------------------------------------------------------------------------

\subsection*{Notations}

% --------------------------------------------------------------------------

Throughout this article, we use the following notations:

\begin{itemize}
	\item For a complex number $ z $, we denote $ \bm{e}(z) \coloneqq e^{2\pi\iu z} $.
	\item We denote $ q $ a complex variable with $ \abs{q} < 1 $.
	\item etc...
\end{itemize}


% --------------------------------------------------------------------------

\section*{Acknowledgement} \label{sec:acknowledgement}

% --------------------------------------------------------------------------


% --------------------------------------------------------------------------

\section{Preliminaries} \label{sec:preliminaries}

% --------------------------------------------------------------------------

Write fundamental results.

% --------------------------------------------------------------------------

\section{} \label{sec:3}

% --------------------------------------------------------------------------



% --------------------------------------------------------------------------

\subsection{} \label{subsec:}

% --------------------------------------------------------------------------



% --------------------------------------------------------------------------

\section{} \label{sec:}

% --------------------------------------------------------------------------


% --------------------------------------------------------------------------

\subsection{} \label{subsec:}

% --------------------------------------------------------------------------



% --------------------------------------------------------------------------

\section{} \label{sec:}

% --------------------------------------------------------------------------


% --------------------------------------------------------------------------

\subsection{} \label{subsec:}

% --------------------------------------------------------------------------



% --------------------------------------------------------------------------

\bibliographystyle{alpha}
\bibliography{myrefs}

% --------------------------------------------------------------------------
\end{document}
% --------------------------------------------------------------------------